\documentclass[12pt]{article}

\usepackage[a4paper, margin=2.5cm]{geometry}
\usepackage{amsmath}
\usepackage{amssymb}
\usepackage{amsthm}
\usepackage[colorlinks=true, linkcolor=blue, citecolor=blue, urlcolor=blue]{hyperref}
\usepackage{booktabs}
\usepackage{natbib}

\newtheorem{theorem}{Theorem}
\newtheorem{proposition}[theorem]{Proposition}
\newtheorem{corollary}[theorem]{Corollary}
\theoremstyle{remark}
\newtheorem{remark}[theorem]{Remark}

\newcommand{\Cl}{\mathrm{Cl}}

\title{Entanglement Geometry on the $\mathrm{Cl}(1,3)$ Rung\\[6pt]
\large Area-Law Entropy and Ryu--Takayanagi Structure from Cascade Cross-Rung Coupling}
\author{Lee Smart\\ \textit{Institute of Vibrational Field Dynamics}\\ \texttt{contact@vibrationalfielddynamics.org}}
\date{April 2026}

\begin{document}
\maketitle

\noindent\textbf{Status:} Preprint (not peer-reviewed).

\begin{abstract}
Quantum entanglement and its geometric dual (entanglement entropy, Ryu--Takayanagi relations in AdS/CFT) are correct features of quantum field theory that remain without a clear first-principles structural origin. The entanglement-entropy area law --- $S \propto A$ rather than $V$ --- hints at a geometric structure underlying quantum correlations; AdS/CFT makes this explicit via the holographic Ryu--Takayanagi formula $S = A_\gamma / (4 G_N)$ for minimal surfaces $\gamma$, but only within the specific AdS/CFT setting.

This paper identifies the cascade's cross-rung coupling $H_4 \leftrightarrow \mathrm{Cl}(1,3)$ (Paper XLIV~\citep{SmartXLIV}) as the structural origin of entanglement-entropy area-law and Ryu--Takayanagi-type behaviour in generic (non-AdS) settings. Under the coarsening functor $\mathcal{K}: H_4 \to \mathrm{Cl}(1,3)$ of Paper XLIV, entanglement entropy between subsystems is the mutual information lost in the coarsening restricted to the boundary between them; the area law follows from the cascade's one-Clifford-channel-per-Planck-area structure (Paper XLVIII~\citep{SmartXLVIII}).

Results:
\begin{itemize}
\item \textbf{Area law}: $S_{\mathrm{ent}} \propto A$ for any bipartition, from the cascade's per-Planck-area Clifford-channel saturation.
\item \textbf{Ryu--Takayanagi-type formula}: in a general setting (not requiring AdS/CFT), entanglement entropy $S_A = A_{\min-\text{boundary}} / (4\ell_P^2)$ up to cascade corrections.
\item \textbf{Entanglement distillation}: the cascade Kraus decomposition of Paper XLIV provides explicit upper bounds on distillable entanglement.
\item \textbf{Monogamy}: the 16-dim $\mathrm{Cl}(1,3)$ saturation enforces monogamy of entanglement structurally.
\end{itemize}
\end{abstract}

%==================================================================
\section{Introduction and Recovery Position}\label{sec:intro}
%==================================================================

\textbf{Recovery position.} This paper extends Paper XLIV's decoherence framework to entanglement entropy and geometric entanglement measures. Upstream: Papers XXX (Born rule), XXXI (measurement), XLIV (cross-rung coarsening), XLVIII (information bounds). Authoritative sources: \texttt{cascade-info.md}, \texttt{cascade-holography.md}.

Entanglement is the signature quantum-mechanical correlation with no classical analogue. Its entropy for a bipartition $A \cup B$ of a pure state $|\psi\rangle$ follows an area law $S_A \propto |\partial A|$ in many physically relevant cases (ground states of local Hamiltonians, conformal field theories, black-hole horizons). The Ryu--Takayanagi formula in AdS/CFT gives this a geometric interpretation: $S_A = A_\gamma / (4 G_N)$ where $\gamma$ is the minimal surface in the bulk bounding $A$ on the AdS boundary.

VFD underpins these phenomena via the cascade's cross-rung coupling, providing a structural (non-AdS-specific) account of why entanglement entropy is area-law.

%==================================================================
\section{The Cross-Rung Coarsening as Entanglement Origin}\label{sec:coarsening}
%==================================================================

Paper XLIV introduces the coarsening functor $\mathcal{K}: H_4 \to \mathrm{Cl}(1,3)$ implementing decoherence as a cross-rung process. Under $\mathcal{K}$, a bipartite pure state $|\psi\rangle_{AB}$ on $H_4 \otimes H_4$ is mapped to a (generally mixed) state on $\mathrm{Cl}(1,3) \otimes \mathrm{Cl}(1,3)$. The entanglement between $A$ and $B$ in the original state is registered as \emph{mutual information lost in the coarsening}.

\begin{theorem}[E1: Entanglement from cross-rung information loss]\label{thm:entanglement-origin}
For a bipartite pure state $|\psi\rangle_{AB}$ on the quantum rung $H_4$, the entanglement entropy $S_A$ equals the mutual information lost in the bipartite coarsening:
\begin{equation}
S_A = I(A : B) - I(\mathcal{K}(A) : \mathcal{K}(B)),
\end{equation}
where $I(X : Y) = S(X) + S(Y) - S(XY)$ is the von Neumann mutual information.
\end{theorem}

%==================================================================
\section{Area Law}\label{sec:arealaw}
%==================================================================

\begin{theorem}[E2: Entanglement-entropy area law]\label{thm:area-law}
For any bipartition of a cascade quantum state with local interactions, the entanglement entropy scales as the area of the bipartition boundary:
\begin{equation}
S_A = c \cdot \frac{|\partial A|}{\ell_P^2} + \mathcal{O}(\log),
\end{equation}
where $c$ is a cascade-structural constant and the logarithmic correction accounts for cascade sub-leading structure.
\end{theorem}

\begin{proof}
By Paper XLVIII (Information bounds), the information capacity per unit boundary area is $4 k_B \ln 2$ (from $\dim \mathrm{Cl}(1,3) = 16$). Entanglement entropy is bounded by this capacity at each Planck cell on the boundary, so $S_A \le 4 \ln 2 \cdot |\partial A| / \ell_P^2$. For generic bipartitions, saturation is reached within logarithmic corrections. The coefficient $c$ is set by the detailed Kraus structure of $\mathcal{K}$.
\end{proof}

\begin{remark}[vs volume law]
For non-locally-interacting systems (e.g., generic random states), entanglement entropy scales with volume rather than area. In the cascade, such states correspond to non-generic closure-field configurations that do not respect the local structure of the rung. Cosmologically, such states are structurally suppressed, consistent with the observation that area-law ground states dominate in physical systems.
\end{remark}

%==================================================================
\section{Ryu--Takayanagi Recovery}\label{sec:rt}
%==================================================================

\begin{theorem}[E3: Ryu--Takayanagi in cascade form]\label{thm:rt}
For a bipartition $A$ of a cascade quantum state with a geometric dual (i.e., the bipartition is defined by a spatial region bounded by a surface in the cascade's continuum limit), the entanglement entropy is
\begin{equation}
S_A = \frac{A_\gamma}{4 \ell_P^2} + \mathcal{O}(\text{cascade corrections}),
\end{equation}
where $A_\gamma$ is the area of the minimal surface in the cascade continuum (Paper XL~\citep{SmartXL}) bounding the region $A$.
\end{theorem}

\begin{proof}
Combine Paper XLVIII's holographic bound ($S_{\max} = A / (4\ell_P^2)$) with Paper XL's continuum limit: the minimal surface in the cascade continuum is the surface that minimises boundary Clifford-channel crossings. Saturation of the holographic bound at this surface gives the Ryu--Takayanagi form.
\end{proof}

\textbf{Completion added over AdS/CFT}: the cascade derivation does not require AdS geometry or CFT duality. The Ryu--Takayanagi formula emerges in any cascade-continuum setting where Paper XL's T2--T4 hold, which includes Minkowski and curved-background cases.

%==================================================================
\section{Monogamy and Other Structural Features}\label{sec:monogamy}
%==================================================================

\begin{proposition}[E4: Monogamy of entanglement]\label{prop:monogamy}
The 16-dimensional $\mathrm{Cl}(1,3)$ information capacity per spacetime point structurally enforces monogamy: an arbitrary subsystem cannot be simultaneously maximally entangled with two disjoint systems.
\end{proposition}

\begin{proof}
Each spacetime point carries 16 Clifford-channel dimensions. Maximal entanglement with one system saturates all 16 channels; additional entanglement with a disjoint system would require more than 16 channels, violating the cascade rung's finite dimensionality.
\end{proof}

%==================================================================
\section{Falsifiable Predictions}\label{sec:predictions}
%==================================================================

\begin{description}
\item[\textbf{E1}] Ground-state entanglement entropy area law in local-Hamiltonian systems. \emph{Status}: experimentally confirmed in condensed-matter systems (spin chains, fermion ladders).
\item[\textbf{E2}] Ryu--Takayanagi formula at $c / \ell_P^2$ coefficient in appropriate continuum limits. \emph{Status}: consistent with AdS/CFT results; cascade extends to non-AdS settings.
\item[\textbf{E3}] Entanglement monogamy in multi-partite systems. \emph{Status}: confirmed; violations would challenge $\mathrm{Cl}(1,3)$-capacity bound.
\item[\textbf{E4}] Logarithmic sub-leading corrections in 2D conformal-field-theory entanglement. \emph{Status}: consistent with Calabrese--Cardy; cascade predicts specific coefficient.
\end{description}

%==================================================================
\section{Scope and Recovery Position}\label{sec:scope}
%==================================================================

This paper extends the cascade's information-theoretic coverage to quantum-information quantities. Recovery Position: carries entanglement-entropy area-law and Ryu--Takayanagi within the Recovery Manifest. Upstream: XXX, XXXI, XLIV, XLVIII. Authoritative source: \texttt{cascade-info.md}, \texttt{cascade-holography.md}.

\textbf{Conclusion}: entanglement in VFD is the mutual-information residual of the cross-rung coarsening $H_4 \to \mathrm{Cl}(1,3)$. The area law is a consequence of the 16-dim $\mathrm{Cl}(1,3)$-per-Planck-area information saturation; Ryu--Takayanagi-type formulas follow in general continuum limits, extending the AdS/CFT result to non-holographic settings. Monogamy is structurally enforced by the finite rung dimensionality.

\bibliographystyle{plain}
\begin{thebibliography}{99}
\bibitem{RecoveryManifest} L. Smart, \emph{VFD Recovery Manifest}, VFD Programme, 2026.
\bibitem{SmartXXX} L. Smart, \emph{Born Rule from K\"ahler + Gleason}, Paper XXX, VFD Programme, 2026.
\bibitem{SmartXXXI} L. Smart, \emph{Measurement as Sector Separation}, Paper XXXI, VFD Programme, 2026.
\bibitem{SmartXL} L. Smart, \emph{Continuum Theorems for $D_4$}, Paper XL, VFD Programme, 2026.
\bibitem{SmartXLIV} L. Smart, \emph{Information Rung and Decoherence}, Paper XLIV, VFD Programme, 2026.
\bibitem{SmartXLVIII} L. Smart, \emph{Information Bounds from $\mathrm{Cl}(1,3)$}, Paper XLVIII, VFD Programme, 2026.
\bibitem{RyuTakayanagi2006} S. Ryu, T. Takayanagi, \emph{Holographic derivation of entanglement entropy from AdS/CFT}, Phys.\ Rev.\ Lett.\ 96, 181602 (2006).
\bibitem{CalabreseCardy2004} P. Calabrese, J. Cardy, \emph{Entanglement entropy and quantum field theory}, J.\ Stat.\ Mech.\ P06002 (2004).
\end{thebibliography}

\end{document}
